\documentclass[11pt,a4paper]{article}
\usepackage[margin=27mm]{geometry}
\usepackage{amsmath,amssymb,amsthm,fontspec,unicode-math,hyperref,booktabs,longtable}
\usepackage{array,calc,fvextra}
\setmainfont{STIXGeneral}
\setmathfont{STIX Two Math}
\setmonofont{Menlo}
\hypersetup{hidelinks}
\setlength{\emergencystretch}{2em}
\providecommand{\tightlist}{\setlength{\itemsep}{0pt}\setlength{\parskip}{0pt}}
\newtheorem{theorem}{Theorem}
\newtheorem{lemma}[theorem]{Lemma}
\newtheorem{corollary}[theorem]{Corollary}
\newtheorem{proposition}[theorem]{Proposition}
\newcommand{\gen}[1]{\langle #1\rangle}
\newcommand{\Sym}{\operatorname{Sym}}
\newcommand{\Aut}{\operatorname{Aut}}
\newcommand{\Fix}{\operatorname{Fix}}

\title{Group representations by recognizable incidence structures}
\author{Human Compression Project}
\date{}
\begin{document}
\maketitle
\begin{abstract}
We recover the finite constructions and free-basis geometries of Leemans,
Stokes and Tranchida through the common restriction theorem, an exact
shadow-kernel theorem and a free-basis normalizer calculation. The dihedral,
symmetric and Platonic examples retain their original attribution and are
checked at the exceptional parameters. For ordered projective quadruples,
we compute the literal obstruction and give a finite-field replacement
whose distinct-length paths retain the ordered endpoints. For the free-group examples, we determine the full type
image and the exact natural normalizer. The finite-pair existence theorem
of D\'iaz Ramos, Molinier and Viruel remains a separate, fully included
application. A complete finite certificate and a uniform doubling
construction refute the unqualified splitting observation in Section 8.
\end{abstract}

\section{Source and common hypotheses}
The source is D. Leemans, K. Stokes and P. Tranchida,
\emph{Groups represented by incidence geometries},
arXiv:2506.10464v1, 12 June 2025, thirteen pages. Its incidence systems
have a finite set of types, but may have infinitely many objects. Write
$C(X)$ for the full correlation group and $K(X)$ for the palette kernel.
A representation of a normal pair $H\lhd G$ is a pair isomorphism
$(G,H)\cong(C(X),K(X))$. A representation of a group $L$ means the pair
$(\Aut(L),\operatorname{Inn}(L))$; a weak representation means compatible
injections of these two groups. The geometric condition that every flag
extends to a chamber is a separate property. In particular, the abstract's
reference to the full group $\operatorname{Out}(L)$ is not the
representation definition used in the body: the outer group is the
quotient $C(X)/K(X)$ when an exact representation is present.

The restriction-and-extension theorem, the complete shadow-kernel theorem
and the unique-shadow corollary of the common paper supply the shared proof. For each finite example
we must recognize the old objects under \emph{all} correlations, reconstruct
the old incidences, prove pointwise rigidity and exhibit every extension.
Counting a chosen group of automorphisms is not by itself a calculation
of the full correlation group. For the infinite examples, Theorems
R2.Geometry, R2.TypeImage and R2.Normalizer of the common paper provide
the complete free-basis argument.

\section{The dihedral family}\label{the-dihedral-family}

Let \(n\ge3\), let \(V=\mathbb Z/n\mathbb Z\), and put
\(U_n=(\mathbb Z/n\mathbb Z)^\times\). For \(k\in U_n/\{\pm1\}\) let \[
E_k=\{\{x,x+k\}:x\in V\}.
\] Each \(E_k\) is a Hamiltonian cycle. If \(n\) is odd, all vertices
have one type. If \(n\) is even, even and odd vertices have separate
types. Each \(E_k\) is a further type. Vertex--edge incidence is
containment, edges of different types are mutually incident, and in the
even case every even vertex is incident to every odd vertex. These are
exactly the incidence systems of {[}L126, §§3.1--3.2{]}.

\textbf{Theorem R1.Dihedral.} For \(n>3\), the full correlation group is
the affine group \[
C(\Gamma_n)=V\rtimes U_n,
\qquad x\longmapsto ax+b.
\tag{D1}
\] Its type-preserving subgroup consists of the affine maps with
\(a=\pm1\), with arbitrary \(b\) for odd \(n\) and even \(b\) for even
\(n\). This gives the exact pair
\((\operatorname{Aut}D_{2n},\operatorname{Inn}D_{2n})\), where the
subscript is the group order. For \(n=3\) it gives the weak
representation stated in the source.

\textbf{Proof.} First the vertices are recognizable. For even \(n\),
their two types have size \(n/2\), while every edge type has size \(n\).
For odd \(n>3\), put \(d=\varphi(n)/2\ge2\). A vertex has degree \(2d\),
and an edge object has degree \(2+n(d-1)\). Their difference is
\((n-2)(d-1)>0\). Thus every correlation preserves \(V\) in both cases.
Each edge object is determined by its two endpoints.

The affine maps in (D1) preserve incidence and send \(E_k\) to
\(E_{ak}\). They act transitively on incident vertex--edge pairs: the
image of the pair \((0,\{0,1\})\) may be prescribed by choosing its
vertex image \(b\) and one of the two endpoint differences \(a\in U_n\).
There are \(n\varphi(n)\) such pairs. A correlation fixing one pair
fixes its edge type, hence acts on its Hamiltonian cycle fixing an edge
and both endpoints. An automorphism of a cycle fixing two adjacent
vertices is the identity, since successive neighbors are forced. It
fixes all vertices and therefore all edge objects. The action on these
pairs is thus free as well as transitive, proving (D1).

An affine map fixes every edge type only if it fixes \(E_1\),
equivalently \(a=\pm1\); these two multipliers indeed fix every \(E_k\).
For even \(n\) it fixes both vertex types exactly when \(b\) is even.
For odd \(n\) there is no further constraint.

Write \(D_{2n}=\langle r,s:r^n=s^2=1,\ srs=r^{-1}\rangle\). Its
automorphisms are exactly \[
r\mapsto r^a,\qquad s\mapsto r^b s,
\quad a\in U_n, b\in V.
\] Their composition is the affine composition law. Conjugation by \(r\)
has parameters \((1,2)\) and conjugation by \(s\) has parameters
\((-1,0)\), so the inner subgroup is precisely the subgroup just
computed. These formulas also verify compatibility of the two pair
isomorphisms.

For \(n=3\) the incidence graph is a hexagon. Its type-preserving group
is \(S_3\) and its full group has order twelve. Since
\(\operatorname{Aut}S_3=\operatorname{Inn}S_3=S_3\), the same inclusion
into the type-preserving subgroup gives compatible injections.
\(\square\)

\textbf{Source qualifications.} The unit group \(U_n\) need not be
cyclic. The source's factor \(C_{\varphi(n)}\) must be replaced by
\(U_n\); for instance \(U_8\cong C_2^2\). The printed equality involving
\(\operatorname{Out}D_6\) in the small-case proof is also unnecessary:
\(\operatorname{Out}S_3\) is trivial. Neither correction changes the
asserted pair in Theorem 3.1.

For some parameters \(\Gamma_n\) is only an incidence system. At
\(n=5\), two differently colored disjoint edges form a flag with no
vertex completing it. At \(n=6\), opposite-parity vertices at difference
three are incident but have no incident edge object in common. This does
not invalidate Theorem 3.1, whose defined notion of representation uses
incidence systems. If a genuine finite geometry is required, the
pair-preserving cleaning and chamber-completion constructions of Díaz
Ramos--Molinier--Viruel {[}F010{]} apply. Their complete proofs are included
in the common paper; no archived construction theorem is left unpaid.

\section{Symmetric groups and the exceptional
degree}\label{symmetric-groups-and-the-exceptional-degree}

For \(n\ge2\) let \(X_n\) be the rank-two incidence geometry of vertices
and edges of the complete graph on \(n\) vertices. The point and edge
classes have sizes \(n\) and \(\binom n2\).

\textbf{Theorem R1.Symmetric.} If \(n\ge4\), then \[
C(X_n)=K(X_n)=S_n.
\] Consequently \(X_n\) is an exact representation of \(S_n\) for
\(n\ge4\), \(n\ne6\). For \(n=3\) it is the preceding weak
representation. The literal \(n=2\) instance is not a representation of
\(S_2\) in the source's sense.

\textbf{Proof.} For \(n\ge4\) the two type sizes differ. Correlations
therefore preserve the vertices. Edges are uniquely recovered from their
endpoint sets, and all vertex permutations extend. This proves the
displayed pair. We retain the classical theorem
\(\operatorname{Aut}S_n=\operatorname{Inn}S_n\) for \(n\ge4\),
\(n\ne6\), also used in {[}L126, Theorem 4.1{]}. For \(n=2\), both
groups of \(X_2\) have order two, whereas both \(\operatorname{Inn}S_2\)
and \(\operatorname{Aut}S_2\) are trivial. Thus the source's parameter
exclusion \(n\ne3,6\) needs a small-parameter qualification. A
one-object rank-one geometry supplies the trivial pair when desired.
If the parameter $n=1$ is included, the declared rank-two model has an
empty edge type and hence fails the source's surjectivity requirement
for the type map; the same one-object rank-one replacement represents
the trivial group.
\(\square\)

We can make the exceptional construction exact, without assuming an
order calculation for \(\operatorname{Aut}S_6\). A \emph{duad} is a
two-element subset of a six-element set, and a \emph{syntheme} is a
partition into three duads. Form \(Q\) with the fifteen duads as points,
fifteen synthemes as lines, and containment as incidence, as in {[}L126,
pp.5--6{]}. Every point lies on a line and every line has three points,
so \(Q\) is a geometry.

\textbf{Theorem R1.Six.} The duad--syntheme geometry is an exact
representation of \(S_6\).

\textbf{Proof.} Distinct duads are collinear exactly when they are
disjoint. Hence a type-preserving automorphism preserves disjointness of
duads. The largest families of pairwise intersecting duads are precisely
the six stars consisting of all five duads through a fixed symbol.
Indeed, if an intersecting family has no common symbol, choose
\(\{a,b\}\) and \(\{a,c\}\) in it and then a member avoiding \(a\). That
member must be \(\{b,c\}\), and any further member must meet all three
edges of this triangle, so lies in the triangle. Such a family has at
most three members. The six five-member stars therefore recover the six
original symbols. Every type-preserving automorphism is the
corresponding permutation of symbols; synthemes are then determined by
their incident duads. Thus \(K(Q)=S_6\).

In the abstract group \(S_6\), the involution classes of transpositions,
products of two disjoint transpositions, and products of three have
sizes fifteen, forty-five, and fifteen respectively. This follows by
choosing the moved symbols and pairing them. Therefore every
automorphism of \(S_6\) permutes the first and third classes as a pair.
A transposition commutes with a triple transposition exactly when it is
one of its three factors: a commuting transposition must exchange two
symbols paired with each other by the triple transposition. Consequently
the duad--syntheme incidence relation is intrinsically the commuting
relation between these two classes. Every automorphism of \(S_6\) acts
on \(Q\) as a correlation, and this action is faithful because
transpositions generate \(S_6\).

Conversely every correlation of \(Q\) normalizes \(K(Q)\) and hence acts
by conjugation on its identified \(S_6\). This conjugation action is
faithful. A type-preserving element in its kernel belongs to the trivial
center of \(S_6\). A type-swapping element in its kernel would induce an
\(S_6\)-equivariant bijection between duads and synthemes. Their
stabilizers cannot be conjugate: a duad stabilizer contains seven
transpositions, whereas a syntheme stabilizer contains only its three
constituent transpositions. Thus no such bijection exists.

The two embeddings are inverse because the first action sends each inner
automorphism to the natural symbol permutation, and conjugation by its
image acts on that subgroup as the original group automorphism. Hence
the pair is exactly \((\operatorname{Aut}S_6,\operatorname{Inn}S_6)\).
\(\square\)

The native finite replay also gives \(|C(Q)|=1440\) and \(|K(Q)|=720\).
These counts corroborate the proof; they are not needed to deduce the
pair isomorphism. We do not rederive the atlas examples for \(A_6\) or
\(A_7\) cited but not constructed in {[}L126{]}. General finite
existence is already supplied by F010, with its original attribution.
The \(A_5\) case appears next with a small explicit construction.

\section{Cube and octahedron}\label{cube-and-octahedron}

Regard the cube graph as the Hamming graph on \(\{0,1\}^3\). Give even
and odd vertices distinct types, and give edges and square faces two
further types. First take the source's literal incidence induced from
the cube: vertex--edge, vertex--face and edge--face incidence are
containment, and different vertex objects are not incident.

\textbf{Theorem R1.Cube.} The literal rank-four incidence system has
pair \[
(C,K)\cong(S_4\times C_2,S_4),
\tag{C1}
\] with the second factor in \(C\) generated by the antipodal map. It
represents \(S_4\times C_2\) exactly. Adding incidence between
opposite-color vertices precisely when they are adjacent cube vertices
gives a geometry with the same pair and twenty-four chambers.

\textbf{Proof.} The type sizes are four, four, twelve and six, so every
correlation preserves the vertex union, the edges, and the faces. The
vertex--edge truncation recovers the cube graph. Edges and faces have
distinct vertex shadows, and every graph automorphism preserves their
families and the parity bipartition. Recognition therefore identifies
the full correlation group with the cube graph group.

Translations of \(\{0,1\}^3\) and coordinate permutations give a group
of order \(8\cdot6=48\). An automorphism fixing zero permutes its three
neighbors. Once these are fixed, each weight-two vertex is the common
neighbor of the appropriate pair other than zero, and the weight-three
vertex is forced. Thus the displayed group is the full graph group. The
parity-preserving subgroup has order twenty-four and acts faithfully on
the four even vertices, hence is \(S_4\). Translation by \((1,1,1)\) is
central, interchanges parity, and supplies the direct factor in (C1).

For \(L=S_4\times C_2\), its center is the second factor and is
characteristic. Every automorphism consequently has the form \[
(s,\epsilon)\longmapsto
(a s a^{-1},\epsilon\,\chi(s)),
\] where \(a\in S_4\) and \(\chi:S_4\to C_2\) is either trivial or the
sign character. The description follows by passing to
\(L/Z(L)\cong S_4\), using completeness of \(S_4\), and then inspecting
the central coordinate. The sign shear is a central involution in
\(\operatorname{Aut}L\), and the other factor is its inner group.
Identifying it with the antipodal map proves compatibility of (C1) with
the required pair.

The literal system has no chamber, since its two vertex types are never
incident. In the repaired system, each chamber consists of an edge, its
two endpoints and one of its two incident faces. There are
\(12\cdot2=24\). Every flag extends: within a square, two
opposite-parity vertices are adjacent; every vertex--edge or edge--face
flag extends through an endpoint or a face; and a vertex--face flag
extends through either boundary edge at that vertex. These observations
also handle singleton and empty flags. The added adjacency is already
recovered by the edge objects, so it changes neither group. \(\square\)

The repair establishes geometry but does not assert firmness: some
panels have a unique completion. The source's group-pair theorem
survives for its literal incidence system even though its use of the
word geometry does not.

\section{\texorpdfstring{Dodecahedron, icosahedron and
\(A_5\)}{Dodecahedron, icosahedron and A\_5}}\label{dodecahedron-icosahedron-and-a_5}

Use the Petersen graph with vertices the two-element subsets of a
five-element set and adjacency given by disjointness. It has ten
vertices and fifteen edges. Its twelve pentagons correspond to
Hamiltonian cycles of the complete graph on the five symbols: the five
edges of such a Hamiltonian cycle are Petersen vertices, and
disjointness puts them in pentagon order. Conversely the five pairs in a
Petersen pentagon use each symbol twice, since pairs through one symbol
form an independent set of size at most two in a pentagon. They
consequently form a Hamiltonian cycle on the five symbols.

The natural \(S_5\) acts transitively on these twelve cycles. A cycle
stabilizer is a dihedral group of order ten, all of whose permutations
on the five symbols are even. Thus \(A_5\) has two orbits \(\mathcal F\)
and \(\mathcal P\), each of size six, and every odd permutation
interchanges them. These are the face and Petrie families in the
source's hemidodecahedral construction.

Take vertex, edge, \(\mathcal F\) and \(\mathcal P\) as four types. Use
containment for their vertex and edge incidences. For a face
\(F\in\mathcal F\) and a Petrie cycle \(P\in\mathcal P\), set \[
F*P\quad\Longleftrightarrow\quad F\text{ and }P\text{ have a common Petersen edge}.
\tag{P1}
\] The rule (P1) is an explicit completion of an incidence left
unspecified in {[}L126, p.7{]}; it is not quoted as a printed rule.

\textbf{Theorem R1.Petersen.} The completed system is a geometry with
120 chambers and group pair \((S_5,A_5)\). It represents both \(A_5\)
and \(A_5\times C_2\) exactly.

\textbf{Proof.} For completeness, the required small cycle calculation
has two outcomes:

\begin{center}
\begin{tabular}{rrr}
\toprule
Pairs in $\mathcal F\times\mathcal P$ & Common vertices & Common edges\\
\midrule
30&3&2\\6&0&0\\\bottomrule
\end{tabular}
\end{center}

Here is the complete six-case calculation with one cycle fixed as
$(1,2,3,4,5)$. We write $ab$ for the Petersen vertex $\{a,b\}$.
Rotations and reversal identify Hamiltonian cycles; the six opposite-family
cycles and their common Petersen vertices are
\[
\begin{array}{c|c@{\qquad}c|c}
12354&12,23,45&12435&12,15,34\\
12543&12,34,45&13245&15,23,45\\
13524&\varnothing&14325&15,23,34.
\end{array}
\]
Each listed ordering has odd sign; the twelve Hamiltonian cycles split
into two six-element orbits because their dihedral stabilizers consist
of even permutations. Thus the list is complete. In each three-vertex
intersection, disjointness gives a path of length two. Transitivity on
the first family gives the table for all thirty-six pairs. The optional
replay enumerates all twelve cycles and verifies the same calculation.

Every Petersen edge belongs to two cycles in each family, by
transitivity and the count \(6\cdot5/15=2\). Each incident face--Petrie
pair has a common edge, and each of their common vertices lies on a
common edge by the table. Hence every flag extends to a
vertex--edge--face--Petrie chamber. The number of chambers is \[
15\cdot2\cdot2\cdot2=120.
\]

The type sizes ten, fifteen, six and six recognize vertices and edges.
The graph group is \(S_5\): its largest independent sets are the five
four-member stars of pairs through a symbol. As in the duad proof, a
pairwise intersecting family with no common symbol has size at most
three. Recovering these stars recovers the five symbols. Faces and
Petrie objects have distinct vertex supports, so restriction is
injective. All of \(S_5\) preserves (P1), while exactly \(A_5\) fixes
both cycle families. The recognition lemma gives the asserted group
pair.

Finally \(\operatorname{Aut}A_5=S_5\) and \(\operatorname{Inn}A_5=A_5\),
a classical retained group fact. For \(L=A_5\times C_2\), the center is
characteristic and the quotient is \(A_5\). There is no nontrivial
homomorphism \(A_5\to C_2\), since \(A_5\) is perfect. Thus
\(\operatorname{Aut}L=\operatorname{Aut}A_5=S_5\), with inner subgroup
\(A_5\). The same pair proves both representations. \(\square\)

Among \(S_5\)-invariant completions of the source's two specified
truncations, (P1) is the only one that gives a geometry. The thirty and
six cross-pairs are the two orbits. To see transitivity on the thirty,
fix one Hamiltonian cycle. Its order-five rotation is transitive on the
five noncomplementary opposite-family cycles: the only Hamiltonian edge
sets fixed by this rotation are that cycle and its complement. Including
the six disjoint pairs produces flags having no common point, while
excluding the thirty intersecting pairs leaves no chambers. This
uniqueness concerns that precise completion problem, not all possible
representations of these groups.

\section{The literal projective construction}
Theorem 6.1 of Leemans--Stokes--Tranchida takes the old objects of
$P=\operatorname{PG}(n-1,K)$ and, for each $\lambda\in K\setminus F$,
adds the ordered collinear quadruples with cross-ratio $\lambda$ as a
new type $Q_\lambda$. An old point is incident with a quadruple exactly
when it is one of its four entries; an old higher-dimensional subspace
is incident exactly when it contains the four points. Any two new
objects of different types are incident. The old shadow therefore
remembers the unordered support, not its ordering.

\begin{theorem}[Exact obstruction over four elements]\label{A126.Projective}
For the literal construction with $n=3$, $F=\mathbb F_2$ and
$K=\mathbb F_4$, the pointwise-old-object kernel in the type-preserving
group is $\Sym(12)^{210}$. Hence the literal system is not a
representation of $\operatorname{PGL}(3,4)$. The same instance refutes
the literal root-type truncation of Corollary 6.2.
\end{theorem}
\begin{proof}
Write $K=\mathbb F_2[a]/(a^2+a+1)$. The old plane has 21 points and
21 lines, each line containing five points. For any four-subset $A$ of
a line $\ell$, all 24 orderings have old shadow $A\cup\{\ell\}$.
Their cross-ratios are neither zero nor one, so they belong to
$\{a,a+1\}$. Interchanging the last two positions inverts the
cross-ratio. Since $a^{-1}=a+1$, this is a bijection between the two
sets of orderings. Each type therefore has exactly twelve objects
with this shadow.

There are $21\binom54=105$ distinct such shadows, at each of two
new types. the complete shadow-kernel theorem of the common paper gives exactly
$\Sym(12)^{210}$ as the pointwise kernel. In particular the
type-preserving group contains $\Sym(12)$, of order 479001600,
whereas
\[
 |\operatorname{PGL}(3,4)|=
 \frac{(4^3-1)(4^3-4)(4^3-4^2)}{4-1}=60480.
\]
Thus even the abstract group orders preclude the proposed isomorphism.
There are $42+2\cdot105\cdot12=2562$ objects in the literal model.

For Corollary 6.2 choose $\alpha=a$. Its minimal polynomial over
$\mathbb F_2$ is $X^2+X+1$, with precisely $a,a+1$ as roots in $K$.
For $n=3$ the displayed truncation keeps both old types and both new
types. It is exactly the same counterexample.
\end{proof}

The source proof's uniqueness assertion is the failed step. The always
valid double interchange
$(p_1,p_2,p_3,p_4)\mapsto(p_2,p_1,p_4,p_3)$ already gives different
objects with the same cross-ratio and old shadow. The finite example
above supplies the complete kernel and an order contradiction; it does
not depend on identifying the full correlation group.

\begin{proposition}[Collapsing twins does not restore the field label]
In this four-element example the image of restriction is all $K(P)$.
The nontrivial field automorphism has a type-preserving lift to the
literal extension.
\end{proposition}
\begin{proof}
At both new types, every four-subset on every line has the same
multiplicity twelve. Every old type-preserving automorphism preserves
those multiplicities, so $H=K(P)$ in that theorem. Explicitly,
apply Frobenius $x\mapsto x^2$ to all coordinates and all quadruple
entries, then interchange the last two entries of each quadruple.
The new cross-ratio is $(\lambda^2)^{-1}=\lambda$ for
$\lambda\in\{a,a+1\}$. The two operations commute and each squares
to the identity; their composite is an involutory type-preserving
lift of Frobenius.
\end{proof}

The ordered-position proposition in the common paper gives a role-preserving ordered-tuple
encoding. By itself it does not give a natural projective representation:
recognizability under all correlations, extension of duality and the palette
action are additional requirements. The distinct-arm construction below
discharges them over finite fields.
For finite groups $L$, including each finite projective group, the
complete finite normal-pair theorem of D\'iaz Ramos--Molinier--Viruel,
proved in the first application, supplies a representation by substituting
$(G,H)=(\Aut(L),\operatorname{Inn}(L))$. This is that theorem's
existence consequence, not a verification of the literal projective model.

Two further scope points are visible in the source definitions. Its
preliminary definition requires a finite palette, whereas $K\setminus F$
can be infinite in Section 6. Also the full semilinear group uses all
field automorphisms; identifying it with $\operatorname{Gal}(K/F)$
requires those automorphisms to fix $F$. The finite $\mathbb F_4/\mathbb F_2$
counterexample satisfies both favorable conditions, so neither
qualification removes the obstruction. The orthogonal, degree-two and
Suzuki variants in Section 6.2 are proposed investigations, explicitly
without the necessary construction details in the source; no natural
model for those suggestions is claimed here.

\begin{proposition}[A finite-palette obstruction over a rational function field]
\label{A126.InfiniteField}
For $n\ge3$, the group $L=\operatorname{PGL}(n,\mathbb Q(t))$ has no
exact incidence representation with finitely many types, even if infinitely
many objects are allowed.
\end{proposition}
\begin{proof}
For each integer $m$, substitution $t\mapsto t+m$ is a field automorphism
and therefore acts entrywise as an automorphism $\psi_m$ of $L$.
These automorphisms have distinct images modulo $\operatorname{Inn}(L)$.
Indeed $\psi_m\psi_k^{-1}=\psi_{m-k}$. If $\psi_h$ were inner, the
projective classes of $D=\operatorname{diag}(t,1,\ldots,1)$ and
$\operatorname{diag}(t+h,1,\ldots,1)$ would be conjugate. Thus one
matrix $\operatorname{diag}(t+h,1,\ldots,1)$ would be conjugate in
$\operatorname{GL}(n,\mathbb Q(t))$ to a scalar multiple $uD$. The repeated eigenvalue has multiplicity
$n-1\ge2$ and the other eigenvalue has multiplicity one, so it forces
$u=1$. The remaining eigenvalue forces $h=0$. Thus
$\operatorname{Out}(L)$ is infinite.

For any incidence system with a finite palette $I$, the homomorphism
from correlations to $\Sym(I)$ has kernel precisely the type-preserving
group. Consequently $C/K$ is finite. An exact pair isomorphism with
$(\Aut(L),\operatorname{Inn}(L))$ would identify this finite quotient
with $\operatorname{Out}(L)$, a contradiction.
\end{proof}
This concerns the definition in the source, whose palette is finite.
It is not an obstruction to a different theory allowing infinitely
many types, nor does it rule out weak representations.

\subsection{A faithful finite-field construction}
Here is a positive replacement of the intended finite-field mechanism.
It changes the attachment and includes both point and dual pencils, so
it is distinct from the literal source construction just refuted.

Let $q=p^f$ be a prime power and $n\ge3$. Take the full proper nonzero
subspace incidence system $P=\operatorname{PG}(n-1,q)$, with dimensions
as types. If $f>1$, choose either
\[
R=\mathbb F_q\setminus\mathbb F_p,
\quad\text{or}\quad
R=\{\alpha,\alpha^p,\ldots,\alpha^{p^{f-1}}\}
\quad\text{where }\mathbb F_p(\alpha)=\mathbb F_q.
\]
For $f=1$ put $R=\varnothing$. For every $\lambda\in R$ take all ordered
quadruples of distinct collinear points with cross-ratio $\lambda$,
and all ordered quadruples of hyperplanes in a pencil with cross-ratio
$\lambda$. The latter is the usual cross-ratio on the corresponding
projective line in the dual vector space. The center color is
$\lambda$ in both cases. Apply the rigid-arms theorem of the common paper with the arm lengths
$2,3,4,5$. Write $Y_{n,q,R}$ for the resulting incidence system.

\begin{theorem}[Finite projective representation with ordered attachments]
\label{A126.OrderedProjective}
For $n\ge3$, restriction identifies the pair of the system above with
\[
(C(Y_{n,q,R}),K(Y_{n,q,R}))
 \cong
(\operatorname{P\Gamma L}(n,q)\rtimes\langle\delta\rangle,
 \operatorname{PGL}(n,q)),
\]
where $\delta$ is the standard point--hyperplane duality of order two.
With the retained classical identification of the full projective-group
automorphism group, this is an exact representation of
$\operatorname{PGL}(n,q)$. A finite incidence geometry with the same pair
then follows from the complete cleaning and chamber theorems of the
common paper.
\end{theorem}
\begin{proof}
If $f=1$, no objects are added, and the assertion is the classical
fundamental theorem for finite projective spaces, including duality.
Suppose $f>1$, so $q\ge4$. Every old dimension type has size
\[
 {n\brack d}_q=
 \prod_{a=0}^{d-1}\frac{q^{n-a}-1}{q^{d-a}-1}
 \equiv1\pmod p.
\]
The denominators are nonzero modulo $p$, which justifies the displayed
congruence of integers. There are $L={n\brack2}_q$ point lines and
also $L$ hyperplane pencils. For a fixed cross-ratio other than zero
and one, a projective line has exactly $(q+1)q(q-1)$ ordered quadruples:
choose the first three distinct points, and the cross-ratio determines
the fourth uniquely. Thus each center color has
\[
 M=2L(q+1)q(q-1)
\]
objects; each internal arm color has $M|R|$ objects. Every new type
size is divisible by $p$. Since a correlation permutes entire type
classes and preserves their sizes, it preserves precisely the union
of old types, characterized by size not divisible by $p$.
The recognition hypothesis of the rigid-arms theorem is proved.

We retain the classical fundamental theorem of projective geometry:
the correlations of the full subspace incidence system are the
semilinear projectivities and their composites with duality. A
semilinear transformation with field part $\sigma$ sends cross-ratios
to $\sigma(\lambda)$. Both choices of $R$ are invariant under every
field automorphism. The duality sends an ordered point quadruple to
an ordered hyperplane-pencil quadruple with the same cross-ratio,
and conversely. This can be seen in coordinates by representing a
pencil by its two-dimensional space of linear forms; the standard
duality uses the same coefficient vectors. Hence every old correlation
preserves the tuple family and permutes its center colors. The rigid-arm
theorem supplies its unique extension and rules out further ones.

The internal colors are fixed by every extension. Fixing the old
colors rules out duality because $n\ge3$ and points and hyperplanes
are different types. The remaining action has a field part $\sigma$;
fixing every center color means fixing $R$ pointwise. If $R$ is the
orbit of $\alpha$, this fixes the field generator $\alpha$, hence
$\sigma=1$. If $R=\mathbb F_q\setminus\mathbb F_p$, fixing it
pointwise also fixes the prime field and therefore all of
$\mathbb F_q$. Consequently the palette kernel is exactly
$\operatorname{PGL}(n,q)$. The converse inclusion follows because
projectivities preserve cross-ratios and dimensions.

Finally we retain the standard finite classical-group identification
$\Aut(\operatorname{PGL}(n,q))=
\operatorname{P\Gamma L}(n,q)\rtimes\langle\delta\rangle$ and
$\operatorname{Inn}(\operatorname{PGL}(n,q))=
\operatorname{PGL}(n,q)$ for $n\ge3$. These are the group-theoretic
inputs identifying the normal pair just proved with the desired pair.
The two common conversion theorems apply to the finite properly colored
graph $Y_{n,q,R}$ and preserve both groups, giving the geometric version.
\end{proof}

The prime-field generation condition in the smaller label set is
intentional. Generating $K$ over a larger prescribed subfield need not
force all automorphisms of $K$ fixing the chosen labels to be trivial.
The theorem asserts the displayed finite-field construction only. It
makes no arbitrary infinite-field representation claim, and it does
not identify its new objects with the source's old quadruple objects.

\section{The source's conjugate
families}\label{the-sources-conjugate-families}

First dispatch the small index sets in source Theorem 7.3. For a single
freely independent generator, the group is infinite cyclic and its sole
omitted-generator subgroup is trivial. Its singleton cosets form a
one-type geometry: every chamber is one object and translation acts
regularly. There is no residue of rank at least two to check. If the
empty index set is allowed, the empty system has the single empty
chamber and trivial group. Thus these endpoints satisfy the source
flag-transitivity assertion; the following nontrivial family uses the
common theory for rank at least two.

Let \(F_n=F(x_1,\ldots,x_n)\) with \(n\ge2\), and put \[
b_{j,\epsilon,i}=x_j^\epsilon x_i x_j^{-\epsilon},
\qquad j\ne i,\quad\epsilon\in\{1,-1\}.
\tag{F9}
\] There are \(r=2n(n-1)\) displayed words. Let \(G\) be the subgroup
they generate.

\textbf{Lemma R2.Core.} The words (F9) freely generate \(G\).

\textbf{Proof.} Construct a finite labeled graph with one central vertex
\(o\) and \(2n\) outer vertices \(v_{j,\epsilon}\). Join \(o\) to
\(v_{j,\epsilon}\) by an oriented edge labeled \(x_j^\epsilon\). At
\(v_{j,\epsilon}\) put a loop labeled \(x_i\) for every \(i\ne j\). Map
this graph to the rose labeled \(x_1,\ldots,x_n\).

At each vertex, all outgoing oriented labels are different: at \(o\)
they are the \(2n\) signed letters; at \(v_{j,\epsilon}\) the edge back
has label \(x_j^{-\epsilon}\), and the loop directions have labels
\(x_i^{\pm1}\) for \(i\ne j\). Thus a reduced edge path maps to a
reduced word. A nonempty reduced closed path cannot map to the identity
in \(F_n\), so the induced homomorphism on fundamental groups is
injective. The spokes form a spanning tree. The associated free basis
consists of a spoke, one outer loop, and the reverse spoke, giving
exactly (F9). Equivalently the graph rank is \[
\bigl(2n+2n(n-1)\bigr)-(1+2n)+1=2n(n-1).
\] This proves the assertion. \(\square\)

The finite graph is the core relevant to the displayed finite generating
set. Attaching loops at every integer vertex of entire coordinate axes,
as a literal reading of the subsequent source description suggests,
would instead give infinitely many independent loops. No such enlarged
graph is used here. Trees may be attached to the finite core to obtain
the subgroup covering without altering its fundamental group.

\textbf{Theorem R2.SourceFamily.} Form the source coset system on \(G\)
with one maximal standard parabolic omitting each word (F9). Then it is
flag-transitive and residually connected, \(G\) acts faithfully and
regularly on chambers, and its full type quotient is \[
\operatorname{Sym}(2n(n-1)).
\tag{F10}
\] The signed permutation group \(C_2^n\rtimes S_n\) of the ambient
basis acts faithfully by correlations. Its action is a specified
subgroup, not the complete type image.

\textbf{Proof.} Lemma R2.Core identifies \(G\) with the free group on
the displayed basis. Theorem R2.Geometry and Theorem R2.TypeImage apply.
An ambient signed permutation \(x_k\mapsto x_{p(k)}^{\delta_k}\) sends
\[
b_{j,\epsilon,i}\longmapsto
b_{p(j),\epsilon\delta_j,p(i)}^{\delta_i}.
\tag{F11}
\] It therefore permutes the omitted-generator parabolics and acts by
(F6). Its palette action is faithful: fixing all indices
\((j,\epsilon,i)\) forces \(p\) to fix every symbol and forces every
\(\delta_j=1\), since for each \(j\) another index \(i\) exists.
\(\square\)

For \(n=2\), the four words are precisely those printed before Theorem
7.5. The specified ambient group acts on their four types as the square
group of order eight. But an arbitrary three-cycle of the four free
generators also acts by a correlation, and such a permutation is not in
that square group. The full quotient in (F10) is \(S_4\) of order
twenty-four. Thus the quotient assertion in Theorem 7.5 is false; it is
not repaired merely by giving a more detailed proof of the action of the
chosen subgroup.

The geometric conclusions and the chosen finite actions of Theorems
7.5--7.6 are recovered. Their weak-representation conclusions also hold:
the displayed subgroup is a finite-rank free group, so Corollary
UR.WeakFree embeds its full automorphism group into the type-preserving
group, compatibly with its inner subgroup. This supplies a different
inner action from the source's chamber-regular construction. The source's
argument from the chosen finite signed action does not itself establish
the weak conclusion, and the natural inner embedding cannot extend by
Corollary R2.NaturalObstruction. Both distinctions are essential; the
literal existence conclusion is not refuted by the natural obstruction.

\section{Thinness does not imply splitting without chamber connectivity}
The last observation in Section 8 of the source is stated for thin,
flag-transitive geometries. Its definition of a geometry requires every
flag to extend to a chamber, but does not require the chamber graph to
be connected. The distinction is necessary.

We use one completely specified finite bipartite graph $B$. Its vertex
set is $\{1,\ldots,896\}$; the complete adjacency lists, together with
explicit permutation generators, are supplied in the mandatory file
\texttt{graph.json}. No graph catalogue or search is a premise.
The complete certificate program checks the following finite facts:
\begin{enumerate}
\item $B$ is simple, connected and cubic, with parts of size $448$ and
$1344$ edges.
\item The displayed permutations preserve adjacency and generate a
subgroup $H$ of order $5376$.
\item Starting with vertex $1$ and its three neighbors in four separate
singleton colors, and all other vertices in a fifth color, repeated
neighbor-color refinement has numbers of cells
\[
5,8,11,14,17,20,26,44,135,398,692,842,890,896.
\]
\item The subgroup $H^+$ preserving the two parts has order $2688$ and
one orbit on the $1344$ unordered edges.
\item The $2688$ elements exchanging the parts have orders
\[
\begin{array}{c|rrr}
\text{order}&4&12&28\\\hline
\text{number}&1024&896&768.
\end{array}
\]
\end{enumerate}
These are exhaustive finite calculations on the supplied graph and
permutations, rather than tests of a sample or of a proposed graph family.
Here is why they prove fullness without relying on a graph-automorphism
package. At each refinement step a vertex receives its previous color
and the multiset of colors of its neighbors. Every automorphism fixing
the four initially distinguished vertices fixes every resulting color.
The final partition is discrete, so the pointwise stabilizer of that
star is trivial. A vertex stabilizer therefore embeds in the symmetric
group on its three neighbors and has order at most $6$. Consequently
\[
 |\Aut(B)|\le 896\cdot6=5376.
\]
The explicitly generated subgroup attains this bound, so $H=\Aut(B)$.
The final two calculations concern the full automorphism group.

\begin{theorem}[A finite thin flag-transitive nonsplit example]
\label{A126.ThinCounterexample}
There is a finite rank-three thin, flag-transitive geometry $\Delta$
with $4480$ objects and $10752$ chambers whose full correlation group
does not split over its type-preserving subgroup. More precisely, for
$N=C_2^{2240}$ and the full groups just certified,
\[
 C(\Delta)=N\rtimes H,\qquad K(\Delta)=N\rtimes H^+,
 \qquad C(\Delta)/K(\Delta)\cong C_2.
\]
The chamber graph has exactly $1344$ connected components.
\end{theorem}
\begin{proof}
Apply the thin-doubling proposition of the common paper to $B$.
Its type sizes are $896,896,2688$, so the edge type is recognized.
There are two objects over each of the $896+1344=2240$ underlying
vertices and edges, and eight chambers over each edge. The proposition
therefore gives all the stated sizes and the exact full group pair.
A complement to $K(\Delta)$ would contain an involution in the
exchanging coset. Its projection to $H$ would be an exchanging
involution, contradicted by the exhaustive order table. This proves
nonsplitting. The connected-case splitting theorem of the common paper
is unaffected: each of the $1344$ edge cubes is one chamber component.
\end{proof}

This refutes the unqualified geometric splitting observation, under the
source's stated definition. It is not a classification of groups admitting
thin exact representations, nor a solution of the broader geometric
representation questions in Section 8. The graph certificate is a required
finite premise of this example. The uniform doubling and connected-case
arguments are written proofs; no finite sample substitutes for them.

\section{Exact scope, attribution and remaining source assertions}
The original constructions, their motivating representation problem and
their proposed further investigations belong to Leemans, Stokes and
Tranchida. The finite normal-pair existence result used here belongs to
D\'iaz Ramos, Molinier and Viruel and is proved in the complete first
application. No source correction is asserted as a new historical
result or as a newly solved question.

The finite examples and free-basis proofs here rework the explicit
phase197 manuscripts. The shadow theorem and its four-element projective
instance reproduce the phase196 argument with its separately recorded
independent review. They are not claimed as mathematics first proved
in this edition. The unrestricted panel-permutation construction additionally recovers the
weak free-group representation conclusions using a different compatible
inner action. This is not inferred from the finite signed subgroup, and
no historical priority is claimed.
The common recognition theorem was already used in
these earlier arguments; the present edition replaces its archived
reference by the complete accepted common paper and retains the full
F010 application. There is no unpaid premise in an old proof bank.

The classical automorphism facts for $S_n$ away from degree six and for
$A_5$ are explicit retained inputs. The positive finite projective repair
also retains the fundamental theorem for the full subspace geometry and
the standard automorphism-group identification for $\operatorname{PGL}(n,q)$,
$n\ge3$. The new ordered attachment and its full rigidity proof are supplied
here; these classical facts do not prove the attachment properties. The exceptional degree-six argument
is given here directly. The cited Atlas constructions for $A_6,A_7$ are
retained historical examples, not reconstructions of their individual
incidence tables. Their existence also follows from the complete finite
normal-pair theorem, without identifying its objects with those Atlas
models. The source's proposed orthogonal and Suzuki variants are not
proved constructions there and remain prospective investigations here.

Source Proposition 7.1 is the elementary coset incidence construction:
right multiplication preserves intersection, and intersecting cosets
have a common representative, proving transitivity on vertices and
incident pairs. The product-intersection and residual-generation criteria
quoted as Theorems 7.2 and 7.4 are retained classical results. For every
use in the free-basis family, the common paper gives direct tree and
residue proofs instead. The interpretation of ``independent'' in
Theorem 7.3 is a freely independent family, as required by its reduced-word
argument; the specific conjugate families satisfy this by the explicit
immersed finite core. No weaker irredundancy assumption is substituted.

The general Artin--Tits statement is a pointer to forthcoming work in
the source and is not a premise of these free-basis proofs. The background
Outer Space realization and maximum finite-subgroup discussion are not
needed to construct the signed permutation subgroup. In particular,
torsion-freeness of the inner subgroup by itself does not say that every
finite outer subgroup lifts to an isomorphic automorphism subgroup.
The finite signed actions proved here make no such inference. The source's
infinite-type and more restrictive finite-geometry questions remain
questions. The common connected-chamber splitting proposition supplies a valid
sufficient hypothesis. The complete thin-doubling construction and finite
certificate refute the unqualified splitting observation; they do not
resolve the remaining representation questions.

The inherited finite programs calculate all groups and flag sets in their
stated bounded models; the free-word program checks 45,846 reduced words.
The uniform proofs above do not infer infinite statements from those
samples. The projective program independently constructs all 2562 objects
and their fibers, but the written counting argument proves the displayed
obstruction without relying on the program's output. The full files are
retained as optional validation. The three basis-permutation Lean
certificates prove the action on free generators, transport of each
complementary generated subgroup and faithfulness. The seven inherited
shadow declarations prove the actual pointwise incidence kernel and its
triviality criterion; they do not prove the image or splitting of
restriction. Neither file formalizes the projective instance, free-basis
geometry, residual connectedness or normalizer. The twelve previously
accepted F010 certificates remain scoped as in the first application.

\begin{thebibliography}{9}
\bibitem{LST} D. Leemans, K. Stokes and P. Tranchida,
\emph{Groups represented by incidence geometries}, arXiv:2506.10464v1,
12 June 2025.
\bibitem{DMV} A. D\'iaz Ramos, R. Molinier and A. Viruel,
\emph{Every finite group is represented by a finite incidence geometry},
arXiv:2511.02410v2, 16 December 2025.
\bibitem{BC} F. Buekenhout and A. M. Cohen,
\emph{Diagram Geometry}, Springer, 2013, Theorem 1.8.10 and Corollary 1.8.13,
as cited in the source.
\bibitem{BDL} F. Buekenhout, M. Dehon and D. Leemans,
\emph{An atlas of residually weakly primitive geometries for small groups},
Acad. Roy. Belg., 1999; retained examples as cited in the source.
\end{thebibliography}

\end{document}
